\section{Exclusión Mutua}

Aunque anteriormente se presentó el concepto de exclusión mutua en el glosario, se presenta una definición más completa. La exclusión mutua es una propiedad y un mecanismo de control de concurrencia diseñado para prevenir que dos o más procesos (o hilos de ejecución) accedan de forma simultánea a un mismo recurso compartido ---como variables globales, archivos, estructuras de datos o dispositivos de hardware--- durante períodos en los que la ejecución simultánea podría provocar inconsistencias, pérdida de datos o condiciones de carrera (\emph{race conditions}). El fragmento de código o intervalo de tiempo en el que un proceso accede a dicho recurso se denomina sección crítica y se abreviará como SC.

\subsection{Requisitos de la exclusión mutua}

Para que un algoritmo o mecanismo de sincronización sea considerado una solución válida, robusta y aceptable al problema de la exclusión mutua, debe cumplir obligatoriamente con los siguientes seis requisitos.

\subparagraph{1. Garantía de Exclusión Mutua}
En un instante dado, solo un proceso de entre todos los que compiten por acceder al mismo recurso u objeto compartido debe tener permiso para entrar y ejecutarse dentro de su SC. 

\subparagraph{2. No Interferencia en la Sección No Crítica}
Si un proceso se detiene, se interrumpe o finaliza mientras se encuentra ejecutando código en su sección no crítica (también llamada sección residual), debe hacerlo de forma totalmente aislada, sin interferir, bloquear ni impedir que los demás procesos continúen con su ejecución normal o intenten acceder a la SC.

\subparagraph{3. Ausencia de Interbloqueo e Inanición}
Ningún proceso que solicite entrar a su SC puede ser demorado de manera indefinida. El algoritmo debe estar libre de interbloqueo e inanición.

\subparagraph{4. Progreso (Acceso sin Demora Inmotivada)}
Si ningún proceso se encuentra dentro de su SC y existe uno o más procesos que solicitan entrar, cualquier proceso que lo solicite debe poder ingresar sin demora ni espera innecesaria causada por el propio mecanismo de control.

\subparagraph{5. Independencia de la Arquitectura y Velocidad}
No se debe realizar ninguna suposición sobre las velocidades relativas de ejecución de los procesos ni sobre el número de procesadores (o núcleos) del sistema. El algoritmo debe ser lógicamente correcto sin importar la frecuencia de reloj, la planificación del sistema operativo o el hardware subyacente.

\subparagraph{6. Tiempo Finito en la Sección Crítica}
Cualquier proceso que ingrese a su SC debe permanecer dentro de ella únicamente por un intervalo de tiempo finito, abandonándola eventualmente para liberar el recurso y permitir que otros procesos en espera puedan acceder a él.

\subsection{Implementaciones por software}

En este apartado se realizará la implementación de cuatro intentos de exclusión mutua por software. De manera resumida, cada intento resolverá un problema pero tendrá otros problemas. Ver esto permite analizar como cumplir los 6 requisitos anteriormente descritos. Titularemos cada intento de la siguiente manera:
\begin{itemize}
  \item Intento 1 - Uso de variable \texttt{turn}
  \item Intento 2 - Banderas con verificación previa (\texttt{flag})
  \item Intento 3 - Banderas con reserva previa
  \item Intento 4 - Retroceso por cortesía
\end{itemize}

El código presentado está realizado usando el lenguaje de programación C y las librerías de hilos del sistema operativo Linux. Para probar las soluciones puede copiar y pegar el código, compilarlo y ejecutarlo usando los siguientes comandos en una terminal.
\begin{bashcode}[title=Compilar y ejecutar ejemplos]
gcc nombre_archivo.c -pthread -O0 -o nombre_archivo 
./nombre_archivo
\end{bashcode}

Adicionalmente, para evitar repetir el mismo código en cada intento, la función \texttt{main} y las importaciones se colocan en el bloque \texttt{Base de los intentos}. Si desea ejecutar el código de un intento particular puede simplemente reemplazar el sector del comentario con el código del intento en cuestión.

\begin{ccode}[title=Base de los intentos]
#include <pthread.h>
#include <stdbool.h>
#include <stdio.h>
#include <unistd.h>

// Código de cada intento aquí

int main() {
  pthread_t t0, t1;
  pthread_create(&t0, NULL, process0, NULL);
  pthread_create(&t1, NULL, process1, NULL);
  pthread_join(t0, NULL);
  pthread_join(t1, NULL);
  return 0;
}
\end{ccode}

\subsubsection{Primer Intento: Alternancia Estricta mediante la Variable \texttt{turn}}

El primer intento se basa en una única variable compartida \texttt{turn}, que indica explícitamente a qué proceso le corresponde el acceso a la sección crítica. El algoritmo impone un esquema de alternancia estricta.

\begin{ccode}[title=Intento I]
volatile int turn = 0; // Variable compartida entre procesos

void* process0(void* arg) {
  // Solo 3 iteraciones a modo de demostración
  for (int i = 0; i < 3; i++) {
    printf("P0 trabaja 3 segundos fuera de la SC\n");
    sleep(3);
    while (turn != 0);
    printf("P0 entra en la sección crítica durante 1 segundo\n");
    sleep(1);
    printf("P0 sale y cambia el turno\n\n");
    turn = 1;
    sleep(1);
  }
  return NULL;
}

void* process1(void* arg) {
  for (int i = 0; i < 3; i++) {
    printf("P1 esta listo y espera para poder entrar a la SC.\n\n");
    while (turn != 1);
    printf("P1 entra en la sección crítica\n");
    sleep(0.2);
    printf("P1 sale y cambia el turno\n\n");
    turn = 0;
  }
  return NULL;
}
\end{ccode}

Si bien esta construcción garantiza la exclusión mutua no cumple con los 6 requisitos de la solución aceptable de exclusión mutua.

Al ejecutar el código del cuadro \texttt{Intento I}, la salida tendrá un aspecto tal como se muestra en el siguiente bloque.

\begingroup
\footnotesize
\begin{verbatim}
---- Salida de la ejecución de Intento I ----
P1 esta listo y espera para poder entrar a la SC.
P0 trabaja 3 segundos fuera de la SC
P0 entra en la sección crítica durante 1 segundo
P0 sale y cambia el turno
P1 entra en la sección crítica
P1 sale y cambia el turno
P1 esta listo y espera para poder entrar a la SC.
P0 trabaja 3 segundos fuera de la SC
P0 entra en la sección crítica durante 1 segundo
 ...
\end{verbatim}
\endgroup 

Al analizar el bloque se pueden observar dos cosas importantes: \texttt{(i)} el proceso P0 es el más lento y tiene el primer turno para entrar a la SC. \texttt{(ii)} el proceso P1 es mucho más rápido que P0, pero si o si debe esperar su turno a pesar de que la SC esté libre.

En consecuencia, un proceso fuera de su sección crítica puede impedir el progreso de otro que desea ingresar, aun cuando el recurso compartido esté libre. Esto viola el segundo y cuarto requisito: no interferencia en la sección residual\footnote{Sección residual: aquella sección de código que está fuera de la sección crítica.} y acceso sin demora.

En cuanto a sus ventajas, el primer intento satisface la exclusión mutua y, en ausencia de fallos, la espera limitada. 

\subsubsection{Segundo Intento: Verificación Previa y Condición de Carrera}

Para solucionar el problema del \texttt{Intento I} y permitir que un proceso más rápido pueda acceder a la sección crítica a su ritmo se reemplaza el sistema de alternancia estricta de turnos por banderas.

El protocolo propuesto consiste en lo siguiente: cada proceso tiene una bandera propia. La bandera tiene dos estados:
\begin{enumerate}
  \item Levantada (true): significa que está dentro de la SC.
  \item Abajo (false): significa que está fuera de la SC o sección residual.
\end{enumerate}
Entonces, se debe verificar primero la bandera del otro proceso y, solo si esta se encuentra en \texttt{false}, entrar a la SC y levantar la bandera propia. De este modo, si un proceso se demora en su sección residual o falla fuera de la sección crítica, el otro no queda bloqueado. Corrigiendo así una de las limitaciones del intento 1.

\begin{ccode}[title=Intento II]
// En este código se ha introducido un usleep() para provocar el error rápidamente sin necesidad de depender del planificador del SO
volatile bool flag[2] = {false, false};

void* process0(void* arg) {
  printf("P0 quiere entrar - Lee bandera de P1: %d\n",flag[1]);
  while (flag[1]); 

  usleep(100000); flag[0] = true; // Levanta bandera
  printf(">>> P0 entra a SC y coloca su bandera en true <<<\n");
  sleep(2);
  printf("<<< P0 sale de SC >>>\n");
  flag[0] = false; // Baja bandera

  return NULL;
}

void* process1(void* arg) {
  printf("P1 quiere entrar - Lee bandera de P0: %d\n",flag[0]);
  while (flag[0]);

  usleep(100000); flag[1] = true;
  printf(">>> P1 entra a SC y coloca su bandera en true <<<\n");
  sleep(2);
  printf("<<< P1 sale de SC >>>\n");
  flag[1] = false;

  return NULL;
}
\end{ccode}

El defecto de esta implementación radica en el orden de las operaciones. La verificación \texttt{while(flag[j])} y el anuncio \texttt{flag[i] = true} constituyen dos operaciones no atómicas. Entre ambas puede intercalarse la ejecución del otro proceso.

La ejecución del código propuesto en el cuadro \texttt{Intento II} demuestra la violación de la exclusión mutua:
\begingroup
\footnotesize
\begin{verbatim}
P0 quiere entrar - Lee bandera de P1: false
P1 quiere entrar - Lee bandera de P0: false
>>> P0 entra a SC y coloca su bandera en true <<<
>>> P1 entra a SC y coloca su bandera en true <<<
<<< P0 sale de SC >>>
<<< P1 sale de SC >>>
\end{verbatim}
\endgroup
Ambos procesos han entrado a la misma sección crítica al mismo tiempo, esto jamás debería ocurrir, viola la primer propiedad: Exclusión Mutua.

Este comportamiento corresponde a una condición de carrera clásica de tipo \textit{check-then-act}. Si bien se resuelve la dependencia de la sección residual, se pierde la propiedad fundamental de exclusión mutua.

\subsubsection{Tercer Intento: Anuncio Previo y Condición de Deadlock}

El segundo intento falla porque la verificación \texttt{while(flag[j])} precede al anuncio \texttt{flag[i] = true}, es decir:
\begin{ccode}
while(flag[otro]);  // Reviso y si hay un Proceso en SC espero
flag[mío]=true;     // Si estoy en la SC levanto mi bandera 
\end{ccode}
Esto permite que ambos procesos lean \texttt{false} simultáneamente. Entonces alguien pensó: ``Sencillo, levantemos primero la bandera'', de modo que se permutan las instrucciones quedando:
\begin{ccode}
flag[mío]=true;   // Anuncio que entraré a la SC levantando la bandera
while(flag[otro]);// Reviso y si hay un Proceso en SC espero
\end{ccode}
El código con este cambio se muestra en el cuadro \texttt{Intento III}.

\begin{ccode}[title=Intento III]
volatile bool flag[2] = {false, false};

void* process0(void* arg) {
  printf("P0 quiere entrar - Coloca bandera en true\n");
  flag[0] = true; usleep(100000);
  printf("P0 lee bandera de P1: %s\n", flag[1] ? "true, espera":"false, continua");
  while (flag[1]); 
  printf(">>> P0 entra a SC <<<\n");
  sleep(2);
  printf("<<< P0 sale de SC >>>\n");
  flag[0] = false;

  return NULL;
}

void* process1(void* arg) {
  printf("P1 quiere entrar - Coloca bandera en true\n");
  flag[1] = true; usleep(100000);
  printf("P1 lee bandera de P0: %s\n", flag[0] ? "true, espera":"false, continua");
  while (flag[0]);
  printf(">>> P1 entra a SC <<<\n");
  sleep(2);
  printf("<<< P1 sale de SC >>>\n");
  flag[1] = false;

  return NULL;
}
\end{ccode}

Con este orden, la exclusión mutua queda garantizada. Considérese el caso en que $P_0$ se ejecuta primero. Tras ejecutar \texttt{flag[0] = true} y comprobar que ningún otro proceso tiene su flag levantada, ingresa a la sección crítica. Si $P_1$ intenta ingresar posteriormente, ejecuta \texttt{flag[1] = true} y queda bloqueado en \texttt{while(flag[0])}, dado que \texttt{flag[0]} permanece en \texttt{true} hasta que $P_0$ abandone su sección crítica. El mismo razonamiento aplica de forma simétrica para $P_1$.

Sin embargo, en este intento la propiedad de exclusión mutua se logra a costa de introducir una condición de bloqueo mutuo (\emph{deadlock}). El fallo se manifiesta cuando ambos procesos levantan su bandera de forma concurrente, antes de que ninguno evalúe la condición de espera. 

El error de este intento se manifiesta en la ejecución del código presentado. La salida que se muestra a continuación demuestra que ambos procesos han quedado a la espera de que el recurso se libere. Mas eso nunca va a suceder ya que verdaderamente no hay ningún proceso en la SC y el recurso está libre.
\begin{verbatim}
P0 quiere entrar - Coloca bandera en true
P1 quiere entrar - Coloca bandera en true
P0 lee bandera de P1: true, espera
P1 lee bandera de P0: true, espera
\end{verbatim}

Este intento cumple los requisitos 1 (exclusión mutua), 2 (No interferencia en sección residual), 4 (progreso) y 5 (Independencia de velocidad) pero viola la tercer condición: Espera limitada. En este caso particular se puede ver que ambos procesos esperan indefinidamente a usar un recurso que realmente está libre.

\subsubsection{Intento IV}

En el cuarto intento ocurre algo similar al intento \texttt{III}. En este caso no se colocará el código completo, solo se introduce el nuevo error, que se denomina \emph{live-lock} (Bloqueo vivo). 

El código de este intento busca solucionar el problema de \emph{deadlock}. Se propone lo siguiente: ``si deseo entrar a la SC, establezco mi bandera como true. Verifico si hay algún otro proceso con su bandera activa. Si lo hay bajo mi bandera, espero un tiempo y vuelvo a intentar''.

Es importante recordar que el \texttt{Intento III} fallaba únicamente cuando dos procesos levantaban su bandera al mismo tiempo antes de poder realizar la comprobación. Eso provocaba que ambos procesos quedaran esperando.

La corrección que se aplica al código del \texttt{Intento III} consiste en agregar la bajada de la bandera y un tiempo de espera al momento de comprobación.
\begin{ccode}[title=Intento IV]
/* En cada proceso se agrega la siguiente lógica 
   Aquí solo se muestra para el proceso P0 */
flag[0] = true;

while(flag[1]) {
  flag[0] = false;   // Cedo el paso
  delay;             // Espero un tiempo
  flag[0] = true;    // Vuelvo a intentarlo
}
\end{ccode}

Y, si bien esta solución es muy astuta, tiene un problema. Vease, que del mismo modo que en todos los intentos anteriores, si ambos procesos intentan entrar a la SC al mismo tiempo, quedan interactuando de forma indefinida como se muestra en la figura \ref{fig:intento4}.

\begin{figure}[!ht]
\centering
\begin{minipage}{0.5\textwidth}
\centering
\begingroup
\footnotesize
\begin{verbatim}
t   P0                P1
---------------------------------
1   flag0=true
---------------------------------
2                     flag1=true
---------------------------------
3   flag0=false
---------------------------------
4                     flag1=false
---------------------------------
5   flag0=true
---------------------------------
6                     flag1=true
---------------------------------
7   flag0=false
---------------------------------
8                     flag1=false
\end{verbatim}
\endgroup
\end{minipage}
\caption{Diagrama de tiempos de procesos del \texttt{Intento IV}.}
\label{fig:intento4}
\end{figure}

En el diagrama de tiempos de la figura \ref{fig:intento4} se aprecia que los procesos intentan entrar, se encuentran una bandera activa y retroceden. Este comportamiento se llama bloqueo activo o bloqueo vivo (\emph{live-lock}) y en este caso no se cumple el quinto requisito de la exclusión mutua ya que depende de asumir que el tiempo de \texttt{delay} sea tal que no ocurra la secuencia de la figura.

Tal vez podría pensar ¿Pero y si se coloca un delay aleatorio no sirve? Nuevamente, esto generaría una probabilidad (aunque pequeña, existente) de que dos procesos acaben ejecutando la secuencia de la figura \ref{fig:intento4}, lo que significa que esta solución no es aceptada para resolver el problema de la exclusión mutua.

\subsubsection{Algoritmo de Peterson}

El algoritmo de Peterson no introduce mecanismos nuevos, sino que combina correctamente dos ideas ya analizadas en los intentos previos:
\begin{itemize}
  \item De los intentos 3 y 4 conserva las \textbf{banderas} (\texttt{flag}): indican la intención de entrar a la sección crítica.
  \item Del intento 1 recupera la idea del \textbf{turno} (\texttt{turn}): se utiliza únicamente para romper el empate cuando ambos procesos quieren entrar simultáneamente.
\end{itemize}
Así, el código resultante se muestra en el cuadro \texttt{Algoritmo de Peterson}.
\begin{ccode}[title=Algoritmo de Peterson]
volatile bool flag[2] = {false, false};
volatile int turn = 0;

void* process0(void* arg) {
    printf("P0 quiere entrar\n");
    flag[0] = true;
    turn = 1;

    while (flag[1] && turn == 1);
    printf(">>> P0 ENTRA <<<\n");
    sleep(2);
    printf("<<< P0 SALE >>>\n\n");
    flag[0] = false;

    return NULL;
}

void* process1(void* arg) {
    printf("P1 quiere entrar\n");
    flag[1] = true;
    turn = 0;

    while (flag[0] && turn == 0);
    printf(">>> P1 ENTRA <<<\n");
    sleep(2);
    printf("<<< P1 SALE >>>\n\n");
    flag[1] = false;

    return NULL;
}
\end{ccode}

En síntesis: \texttt{flag} indica quién quiere entrar y \texttt{turn} decide quién entra primero si ambos quieren hacerlo al mismo tiempo.

Para comprender la idea del algoritmo suponga el caso del proceso $P_0$. Para entrar a la sección crítica realiza tres operaciones:
\begin{ccode}[]
flag[0] = true; // 1. Anuncia intención de entrar
turn = 1; // 2. Cede la prioridad al otro proceso
while (flag[1] && turn == 1); // 3. Espera si corresponde
\end{ccode}

La asignación \texttt{turn = 1} cede explícitamente la prioridad a $P_1$. La condición del \texttt{while} implica que $P_0$ espera únicamente si se cumplen \textbf{ambas} condiciones: el otro proceso quiere entrar (\texttt{flag[1] == true}) \textit{y} además le corresponde el turno (\texttt{turn == 1}).

En cada caso, ocurriría lo siguiente:
\begin{enumerate}
    \item Si solo $P_0$ quiere entrar: \texttt{flag[1] == false}, por lo que la condición \texttt{while(flag[1] \&\& turn == 1)} es \texttt{false}. $P_0$ entra inmediatamente.

    \item Si $P_1$ está dentro de la sección crítica y $P_0$ quiere entrar: \texttt{flag[1] == true} y \texttt{turn == 1} es \texttt{true}, por lo que $P_0$ permanece en espera activa hasta que $P_1$ ejecute \texttt{flag[1] = false} al salir.

    \item Y el caso importamte, ambos quieren entrar simultáneamente: Inicialmente \texttt{flag[0] = flag[1] = false}.
      \begin{bashcode}
P0: flag[0]=true; turn=1;
P1: flag[1]=true; turn=0;
      \end{bashcode}
  El valor final de \texttt{turn} dependerá de quién escriba último. Si se supone que el último valor es \texttt{turn = 0}, entonces $P_0$ evalúa:
\begin{ccode}
  while(flag[1] && turn==1); // es falso porque turn = 0, P0 entra
\end{ccode}
  Mientras que $P_1$ evalúa \texttt{while(flag[0] \&\& turn==0)} como verdadero y espera. Al salir, $P_0$ hace \texttt{flag[0]=false}, liberando a $P_1$.
\end{enumerate}

De esta forma, el algoritmo garantiza exclusión mutua, progreso y espera acotada, resolviendo el problema sin requerir instrucciones de hardware especiales.

